\subsection{Настоящий Пионер}

Всё, что отлажено в симуляторе, переносится на дрон. Но дрон --- электромеханическое устройство с вращающимися пропеллерами, поэтому порядок работы с ним строже.

\subsubsection{Техника безопасности}

\begin{warnbox}[Критически важно]
Пренебрежение этими правилами приводит к травмам и поломке оборудования. Они не обсуждаются.
\end{warnbox}

\begin{itemize}
    \item \textbf{Механика --- только без питания.} Любые монтажные работы, замена деталей, подключение модулей --- на обесточенном аппарате.
    \item \textbf{Отладка --- без пропеллеров.} При загрузке и проверке скриптов, тестировании моторов и настройке пульта пропеллеры должны быть сняты. Случайный запуск моторов с пропеллерами --- неконтролируемый взлёт.
    \item \textbf{Аккумуляторы.} Только исправные LiPo-батареи, зарядка под присмотром. Вздувшийся или повреждённый аккумулятор не используется.
    \item \textbf{Порядок включения.} Сначала пульт, потом питание дрона. Выключение --- в обратном порядке.
    \item \textbf{Рабочее место.} На столе нет лишних предметов и проводов; с пропеллерами дрон создаёт сильный поток воздуха.
\end{itemize}

\subsubsection{Pioneer Station}

Основная среда для Lua-программ на настоящем дроне --- \textbf{Pioneer Station}: редактор с подсветкой синтаксиса, загрузчик скриптов на борт, консоль для \texttt{print} и телеметрии, менеджер прошивок. Дистрибутив бесплатно доступен на сайте поддержки Геоскан.

Порядок подключения:
\begin{enumerate}
    \item Убедитесь, что пропеллеры сняты (если вы не планируете немедленный вылет на полигоне).
    \item Подключите аккумулятор. Прозвучит мелодия инициализации, светодиоды начнут мигать.
    \item Соедините дрон с компьютером кабелем micro-USB.
    \item Запустите Pioneer Station и выберите COM-порт (обычно определяется сам).
    \item Вставьте программу в редактор и нажмите «Загрузить». Скрипт передаётся в память автопилота и запускается; результат --- в окне «Консоль».
\end{enumerate}

\subsubsection{Python на настоящем дроне}

Python-программа работает не на борту, а на вашем компьютере, и общается с дроном по Wi-Fi через библиотеку \texttt{pioneer\_sdk}. Установка и подключение разобраны в главе~\ref{sec:python-sdk}. Ту же программу можно направить и на симулятор --- тогда по Wi-Fi полетит виртуальный дрон на сайте (раздел~\ref{sec:python-external}).

\subsubsection{Изучение чистого Lua}

Синтаксис Lua можно тренировать и без дрона: симулятор для этого подходит (терминал показывает \texttt{print}), а на компьютере удобны ZeroBrane Studio, Visual Studio Code с расширением Lua или онлайн-интерпретаторы. Код оттуда для дрона всё равно переносится в симулятор или Pioneer Station.

\subsection*{Итоги}
\begin{itemize}
    \item Программа для Пионера одна, исполнителей два: симулятор в браузере и настоящий дрон.
    \item В симуляторе: редактор $\to$ язык $\to$ \ui{Запустить}; вывод --- в терминале, ошибки --- в уведомлениях.
    \item На дроне: безопасность, Pioneer Station для Lua, \texttt{pioneer\_sdk} для Python.
\end{itemize}
